\scanpage{269}
\mspage{10}

2. \deleted{\uline{\unsure{Composition} des intervalles orientés}} \inserted{\uline{Sous-catégories pleines d'un intervalle}}

Soit $I$ une \inserted{catégorie} intervalle orientée, \inserted{de longueur $n\geq1$. Nous identifions $\operatorname{Ob}I$ à l'ens d'entiers $[0,n]$.} Pour une chaîne \deleted{de} \inserted{$\underline{u}_i,\underline{u}_{i+1},\ldots,\underline{u}_{i+l-1}$} \unsure{flèches} de transition de $I$, on appelle \uline{domaine} de la chaîne, l'ens des sources et des buts de ces éléments. C'est donc un intervalle d'entiers, $[i-1,i+l-1]$, de longueur \inserted{$l$} égale à celle de la chaîne, dont l'origine est déterminée par la première flèche $u_i$ de la chaîne (c'est sa source, si $u_i$ est directe, et son but si $u_i$ est rétrograde), et dont l'extrémité est déterminée par la dernière $u_{i+l-1}$.

Les intervalles d'entiers \inserted{(vus comme des sommets de $I$)} \unsure{correspondant|correspondants} aux chaînes structurales sont appelés les \uline{intervalles structuraux} \inserted{ou \uline{intervalles critiques}} de $I$. Contrairement aux chaînes structurales, ces intervalles ne forment pas une partition de $I$, puisque deux intervalles structuraux peuvent avoir une extrémité commune. Si les chaînes structurales sont

\[C_1,C_2\ldots,C_r\qquad (r=\text{long. utile de }I)\]

de longueurs

\[l_1,l_2,\ldots,l_r,\]

alors les intervalles structuraux sont

\[\underbrace{[0,l_1]}_{I_1}\underbrace{[l_1,l_1+l_2]}_{I_2}\unsure{,}\ldots\ \underbrace{[l_1+\cdots+l_{r-1},l_1+\cdots+l_{r-1}+l_r]}_{I_r},\]

une suite d'intervalles d'entiers dont deux consécutifs ont une extrémité commune, et tels que, si $i<j$, chaque élément de $\unsure{I|I_i}$ soit $\leq$ chaque élément de $I_j$. On appelle \uline{sommets critiques} de $I$ les extrémités des intervalles structuraux.

\uline{Prop. \deleted{\unsure{2}}3} Soit $s\in I$, \inserted{soit $I'=\deleted{\unsure{{}_{s}\backslash}}I_{/s}$ (sous-cat. pl. de $I$ des $s'$ tels que existe $s'\longrightarrow s$), $I''=I_{/s}$, $I'_0=I_{/s}\backslash\{s\}$, $I''_0=\unsure{{}_{s}\backslash I|I_{/s}}\backslash\{s\}$.} \deleted{Pour que $s$ soit critique, il f. et s. que les deux sous-catégories} \deleted{Pour que} Pour que $s$ soit critique, il f. et s. que $I\backslash\{s\}$ soit cat. somme des $I'_0$ et $I''_0$ ; \deleted{et que celle au moins \unsure{un vide}}

% slanted
\margin{i.e. qu'on ait : \unsure{a)} Il n'y a pas de flèche entre un élément de $I'_0$ et un élément de $I''_0$. \deleted{b) $I'_0$ et $I''_0$ sont \unsure{pleines} \illegible{1}} (\unsure{évidemment} \illegible{1} par \unsure{la nécessité}.)}

% archive sheet number: 268
